% Figure polypeps_birth_circuit.
% Owner: none (note-only figure).
% Used by: gapped-2d-notes PP Sec. 7 (Lemma 6, homogeneous birth).
% Source: arXiv OpenAI2026PolynomialPEPS, sections/05-frames.tex:396-447
% (Lemma 6.5 and its proof, eq. (6.13) birth-map; Lemma 6.4, eq. (6.11),
% 05-frames.tex:352-370).
% Formula: G_b = Vhat^dagger (|s><s|_{T B_T} (x) id_{E B_E}) Vhat, Vhat = id_{TE} (x) V,
% with V : H_U -> B_T (x) B_E the isometry and s in H_T (x) B_T the normalized
% vector of the splitting lemma; ||(G_b - 1) Omega|| <= L^{-30} when I(T:E) <= L^{-60}.
% Ink-to-index: wires, top to bottom, carry H_T, (blank spacer row), B_T, B_E,
% H_U, H_E. The spacer row keeps the two party contours at least half a pitch
% apart. Invisible (skin=none) vertices only extend wires to the panel edges
% and hold the ends of the cap and the cup. Right to left: V takes U to
% (B_T, B_E); the cap <s| consumes T and B_T; the cup s prepares T and B_T;
% V^dagger takes (B_T, B_E) back to U. The E wire is untouched. All operator
% names are inscribed (wide=2 boxes), and s, <s| sit inside their bowls. Every
% wire is black: owners are party contours, not wire hues (H2, G2).
% Parties: two dashed red contours (closed strings, drawn like the party
% glyph). P_o holds V, the effect <s|, the cup s with its B_T leg, V^dagger,
% the input T and the U wire. Q_o, the new owner, holds the output T stub.
% P_o runs between the two legs of the cup and crosses its T leg just west of
% the arc, so s is the one pair source and its T index leaves P_o for Q_o;
% everything else is private at P_o. E belongs to other parties.
% Every contour crossing (input T, input U, output U, output T) is declared an
% under-crossing, so the build raises no TKZ-CROSS-TOUCHALL warning.
% Trailing joiner "= G_b": the circuit is the canonical birth map itself, eq.
% (6.13), which acts as the identity on E (Vhat = id_{TE} (x) V); a second
% panel drawing a G_b box would not fit the column at the 6 mm pitch.
% Boundary signature: (phys:e) x3 in (T, U, E) and (phys:w) x3 out (T, U, E);
% B_T and B_E are internal. Every open leg is at least half a pitch long.
% The death is the mirror image (not drawn): <s| becomes a bipartite pair
% effect held by Q_o and P_o, and s is reinserted wholly at P_o.
% Related: polypeps_buffer_splitting draws the splitting lemma V Omega ~ s (x) s'
% that this circuit uses; polypeps_allowed_monomial draws a general monomial.
\begin{tenkzequation}
    \tenkzkernel
    \begin{tenkz}[rows={wire,wire,wire,wire,wire,wire}, cols=10, bonds=none, size=l]
        \tn[at=(3,2), name=bcVdag, skin=box, wide=2, wires=3,
            ports={0@1:physical, 0@2:physical, 180@3:physical}]{V^\dagger}
        \tn[at=(3,8), name=bcV, skin=box, wide=2, wires=3,
            ports={0@3:physical, 180@1:physical, 180@2:physical}]{V}
        \tn[at=(1,4), name=bcSrcT, skin=none, ports={180:physical, 0:physical}]{}
        \tn[at=(3,4), name=bcSrcB, skin=none, ports={180:physical, 0:physical}]{}
        \tn[at=(1,7), name=bcEffT, skin=none, ports={180:physical, 0:physical}]{}
        \tn[at=(3,7), name=bcEffB, skin=none, ports={180:physical, 0:physical}]{}
        \tn[at=(1,1), name=bcWt, skin=none, ports={180:physical, 0:physical}]{}
        \tn[at=(5,1), name=bcWu, skin=none, ports={180:physical, 0:physical}]{}
        \tn[at=(6,1), name=bcWe, skin=none, ports={180:physical, 0:physical}]{}
        \tn[at=(1,10), name=bcEt, skin=none, ports={180:physical, 0:physical}]{}
        \tn[at=(5,10), name=bcEu, skin=none, ports={180:physical, 0:physical}]{}
        \tn[at=(6,10), name=bcEe, skin=none, ports={180:physical, 0:physical}]{}
        \tnwire{bcEt.0}{open e}
        \tnwire{bcEu.0}{open e}
        \tnwire{bcEe.0}{open e}
        \tnwire[name=bcInT]{bcEt.180}{bcEffT.0}
        \tnwire[name=bcInU]{bcEu.180}{bcV.0@3}
        \tnwire[name=bcBt]{bcV.180@1}{bcEffB.0}
        \tnwire[route=arc, name=bcEffect]{bcEffT.180}{bcEffB.180}
        \tnwire[route=arc, name=bcSource]{bcSrcT.0}{bcSrcB.0}
        \tnwire[name=bcBe]{bcV.180@2}{bcVdag.0@2}
        \tnwire{bcSrcB.180}{bcVdag.0@1}
        \tnwire[name=bcOutT]{bcSrcT.180}{bcWt.0}
        \tnwire[name=bcOutU]{bcVdag.180@3}{bcWu.0}
        \tnwire{bcEe.180}{bcWe.0}
        \tnwire[name=bcT]{bcWt.180}{open w}
        \tnwire[name=bcU]{bcWu.180}{open w}
        \tnwire[name=bcE]{bcWe.180}{open w}
        \tnmark[form=label, label pos=w]{on bcT 1}{$T$}
        \tnmark[form=label, label pos=w]{on bcU 1}{$U$}
        \tnmark[form=label, label pos=w]{on bcE 1}{$E$}
        \tnmark[form=label, label pos=s]{on bcBe 0.5}{$\mathcal B_E$}
        \tnmark[form=label, label pos=e]{0.4 s of 0.4 w of (3,4)}{$\mathcal B_T$}
        \tnmark[form=label, label pos=e]{on bcEffect 0.5}{$\langle s\rvert$}
        \tnmark[form=label, label pos=w]{on bcSource 0.5}{$s$}
        \tnwire[kind=string, closed, route=orth, stroke=dashed, species=selected,
            name=bcParty,
            cross={under at crossing of bcParty and bcInT,
                    under at crossing of bcParty and bcInU,
                    under at crossing of bcParty and bcOutU,
                    under at crossing of bcParty and bcOutT},
            via={0.5 n of 0.4 w of (1,4), 0.5 n of 0.5 e of (1,9),
                    0.5 s of 0.5 e of (5,9), 0.5 s of 0.7 w of (5,2),
                    0.7 w of (2,2), 0.4 w of (2,4)}]
        \tnmark[form=label, label pos=n]{0.5 n of 0.5 e of (1,7)}{$P_\circ$}
        \tnwire[kind=string, closed, route=orth, stroke=dashed, species=selected,
            name=bcPartyQ,
            cross={under at crossing of bcPartyQ and bcOutT},
            via={0.4 n of 0.5 w of (1,2), 0.4 n of 0.1 w of (1,3),
                    0.4 s of 0.1 w of (1,3), 0.4 s of 0.5 w of (1,2)}]
        \tnmark[form=label, label pos=n]{0.4 n of 0.2 e of (1,2)}{$Q_\circ$}
    \end{tenkz}
    $=\;G_{\mathrm{b}}$
\end{tenkzequation}
